\section{Discussion}
\label{sec:discussion}

\paragraph{Three behavioral modes.} Resisters (Claude Haiku, Claude Sonnet, GPT-5.4) show negative per-turn slopes on the two easier benchmarks: early flips often reverse. Capitulators (GPT-5.4 Mini, Nano) show the opposite pattern, each turn raises flip probability further. Fragmenters lose a coherent majority answer rather than settling on either the correct or planted one. These are not scale artifacts, resisters and capitulators sit at adjacent capability tiers within the same family, but the split collapses on HLE, where five of six targets fragment on most $T_0$-correct queries. We read this as evidence within our six targets, not a claim about frontier models generally. A single-turn benchmark cannot separate these modes, since all three can share the same $T_1$ flip rate while diverging over later turns.

\paragraph{Uncertainty predicts susceptibility, and pressure appears to amplify it.} The RQ1 contrast, uncertain questions flip roughly twice as often as certain ones, is robust across models and benchmarks. Whether escalating pressure specifically amplifies that pre-existing uncertainty is weaker but present: significant for only one model in isolation, but the pooled estimate across all five conditions is significant and survives dropping the noisiest one. The practical implication does not depend on fully resolving this: a benchmark that fixes a question set rather than sampling across the uncertainty distribution will overstate susceptibility on questions a model is already unsure of, so controlling for baseline uncertainty matters regardless.

\paragraph{Miscalibration under pressure.} The correctness/confidence dissociation in RQ4 extends \citet{mei2025reasoninguncertainty}'s finding that wrong answers can carry high self-verbalized confidence to the social-pressure setting: the useful signal is the gap between belief and confidence, not confidence alone. A confidence-threshold abstention rule would catch almost none of the flips we identify.
